%! Characteristics of Rational Functions — Unit 4, Lesson 4.0 — Warm-Up (Answer Key)
\documentclass[10pt]{article}
\usepackage{algebra2-article}
\usepackage{algebra2-key}   % pulls in -boxes; do NOT also load -boxes
\newcommand{\TallMath}[1]{$\displaystyle #1\rule[-1.4em]{0pt}{3.2em}$}

\begin{document}

\pageheader{Unit 4, Lesson 4.0}{Warm-Up --- Answer Key}
\namedateperiod

\vspace{0.10in}

\begin{notesbox}{Getting Ready: Skills We'll Use Today}
\small These three skills power today's lesson on reading rational graphs. Work quickly.

\vspace{4pt}
\textbf{1. Where is a denominator zero?} A fraction $\dfrac{P}{Q}$ has \emph{no value} when
$Q=0$. Set each denominator equal to zero and solve --- factor first when you can.
\begin{enumerate}[label=(\alph*), leftmargin=2.1em, itemsep=3pt]
  \item $x-1=0$ \quad$\Rightarrow$\quad $x=$ \ans{1}
  \item $x^2-9=0$ \quad factored: \ans{$(x-3)(x+3)$} \quad$\Rightarrow$\quad $x=$ \ans{$3,\ -3$}
  \item $x^2+x-6=0$ \quad factored: \ans{$(x+3)(x-2)$} \quad$\Rightarrow$\quad $x=$ \ans{$-3,\ 2$}
\end{enumerate}

\vspace{4pt}
\textbf{2. What happens \emph{near} a forbidden input --- and \emph{far} from it?} Let
$y=\dfrac{1}{x-1}$, which is undefined at $x=1$. Fill in both tables.
\begin{center}
\renewcommand{\arraystretch}{1.3}
\begin{tabular}{|c|c|c|c|c|}
\hline
\multicolumn{5}{|c|}{\small \emph{closing in on} $x=1$} \\ \hline
$x$ & $0.9$ & $0.99$ & $1.01$ & $1.1$ \\ \hline
$y$ & \ans{-10} & \ans{-100} & \ans{100} & \ans{10} \\ \hline
\end{tabular}
\hspace{0.25in}
\begin{tabular}{|c|c|c|c|}
\hline
\multicolumn{4}{|c|}{\small \emph{running off to the right}} \\ \hline
$x$ & $11$ & $101$ & $1001$ \\ \hline
$y$ & \ans{0.1} & \ans{0.01} & \ans{0.001} \\ \hline
\end{tabular}
\end{center}
As $x$ closes in on $1$, the outputs grow \ans{without bound (huge $\pm$ values)}. \quad As $x$ runs
off to the right, the outputs close in on $y=$ \ans{0}.

\vspace{4pt}
\textbf{3. Factor and cancel --- then check the \emph{original}.} Simplify
$\dfrac{x^2-4}{x-2}$.
\begin{enumerate}[label=(\alph*), leftmargin=2.1em, itemsep=3pt]
  \item Factor the numerator: $x^2-4=$ \ans{$(x-2)(x+2)$}, \; so after cancelling the expression is
        \ans{$x+2$}.
  \item Substitute $x=2$ into the \textbf{original} fraction. What do you get?
        \ans{$\tfrac{0}{0}$ --- undefined}
  \item So the simplified form and the original agree at every input \emph{except}
        $x=$ \ans{2}.
\end{enumerate}

\end{notesbox}

\begin{teachernote}
Each item seeds one of today's three new ideas, so debrief all three aloud.
\textbf{Item 1} is the domain-restriction engine: every excluded value of a rational function comes
from solving \emph{denominator} $=0$ (the 3.2 factoring toolkit, reused). \textbf{Item 2} splits the
two asymptote behaviors --- the left table blows up (\emph{vertical} asymptote at $x=1$) while the
right table settles down (\emph{horizontal} asymptote at $y=0$); the sign flip from $-100$ to $100$
across $x=1$ is exactly why the graph will come in \emph{two pieces}. \textbf{Item 3} previews
\emph{holes}: cancelling is legal, but $x=2$ was already thrown out of the domain and cannot come
back --- the graph is the line $y=x+2$ with one point punched out. Do \emph{not} name ``vertical
asymptote,'' ``horizontal asymptote,'' or ``hole'' yet; let the Hook graph name them.
\end{teachernote}

\end{document}
